Los resultados experimentales muestran que una mejor complejidad no 
necesariamente indica un menor tiempo de ejecución. En ordenamiento, \texttt{std::sort}
presentó en general los menores tiempos, mientras que QuickSort
y PatienceSort superaron el límite de 30 segundos en algunas
entradas de tamaño $10^7$. Esto muestra que no basta con considerar
la cantidad de elementos, pues también influyen su distribución
y el orden inicial. En QuickSort, elegir un pivote central no evitó
las particiones desbalanceadas cuando había muchos valores repetidos.
Por su parte, PatienceSort mostró dificultades con las entradas
ascendentes del dominio D7, donde la construcción de muchas pilas
aumenta los costos de la implementación.

En multiplicacion de matrices, Naive se desenvolviò considerablemente más
rápido que Strassen para todos los casos experimentados. Para
$n=1024$, Naive terminó las pruebas en aproximadamente $0.872$ s
en promedio, mientras que Strassen superó el límite de 60 segundos
en todas sus ejecuciones. Aunque Strassen reduce la cantidad de
multiplicaciones recursivas, las copias de matrices, las operaciones
auxiliares y la recursión hasta matrices de tamaño 1 generan un
costo que resultó importante en estas pruebas. Por lo tanto,
estos resultados no permiten concluir que Naive siempre será
mejor, sino que fue más conveniente con las implementaciones
y tamaños utilizados.

En memoria también se observaron diferencias: Strassen presentó
un mayor consumo que Naive en los tamaños comparables, mientras
que en ordenamiento el uso de arreglos auxiliares y pilas influyó
en los resultados. Sin embargo, se midió la memoria máxima del
proceso completo, por lo que no todo ese consumo corresponde
exclusivamente al algoritmo. Además, los casos con
\texttt{TIMEOUT} permiten reconocer una limitación práctica,
pero no conocer su tiempo final ni comparar su memoria.

En conjunto, pudimos observar que la complejidad teórica ayuda
a entender el crecimiento esperado, pero no describe por sí sola
todo lo que ocurre al ejecutar un programa. La elección de un
algoritmo también debe considerar cómo está implementado, qué
entradas recibe y cuánta memoria necesita. Como continuación,
sería interesante evaluar una partición de QuickSort que agrupe
los valores iguales y una versión de Strassen que utilice Naive
para las matrices pequeñas, para comprobar cuánto cambian los
resultados al reducir estos costos.